%%%%%%%%%%%%%%%%%%%%%%%%%%%%%%%%%%%%%%%%%%%%%%%%%%%%%%%%%%%%%%%%
%%%%%%%%%%%%%%%%%%%%%%%%%%%%%%%%%%%%%%%%%%%%%%%%%%%%%%% Oversigt
%\section{Introduktion}
%\begin{frame}
%  \vspace{25mm}
%  \begin{center}
%    \Huge{Part 0:\\Introduktion}
%  \end{center}
%\end{frame}

\section{Arbejdsmodel}
\begin{frame}
  \frametitle{Arbejdsmodel}
  Hvordan kommer det til at foregå?
  
  Vi har forberedt noet fagligt materiale til jer.
  
  Det er udformet som moduler der hver indeholder lidt materiale og nogle øvelser.
  
  I sidder i grupper af to og arbejder sammen om at løse disse opgaver.
  
  \ldots eller også snakker I med en repræsentant for uddannelserne.
\end{frame}

\section{Livebook}
\begin{frame}
    \frametitle{Livebook}
    \vspace{5mm}
    Vi kommer til at arbejde igennem et værktøj der hedder Livebook, som I gerne allerede skulle have installeret.
    
    \vspace{7mm}
    Ellers er det her:
    \begin{center}
      \url{https://livebook.dev/\#install}
    \end{center}
    
    \vspace{7mm}
    Start her:
    \begin{center}
      \scalebox{0.9}{\url{https://github.com/aslakjohansen/sdu-sest-student-visit/tree/main/src/livebook}}
    \end{center}
\end{frame}

\section{Repositorie}
\begin{frame}
    \frametitle{Repositorie}
    \begin{center}
      \scalebox{0.8}{\includeSVG{repo-barcode}}
      \url{https://github.com/aslakjohansen/sdu-sest-student-visit}
    \end{center}
\end{frame}

%%%%%%%%%%%%%%%%%%%%%%%%%%%%%%%%%%%%%%%%%%%%%%%%%%%%%%%%%%%%%%%%
%%%%%%%%%%%%%%%%%%%%%%%%%%%%%%%%%%%%%%%%%%%%%%%%%%%%%%% Oversigt
%\section{Introduktion}
%\begin{frame}
%    \vspace{25mm}
%    \begin{center}
%        \Huge{Part 0:\\Introduktion}
%    \end{center}
%\end{frame}

%\subsection{Overblik}
%\begin{frame}
%    \frametitle{Overblik}
%    \begin{center}
%      \scalebox{1.6}{\includeSVG{pubsub_flow}}
%    \end{center}
%    \vspace{8mm}
%    Vi skal lave et \textsl{dashboard}!
%\end{frame}

%\subsection{Repositorie}
%\begin{frame}
%    \frametitle{Repositorie}
%    \begin{center}
%      \scalebox{0.8}{\includeSVG{repo-barcode}}
%      \url{https://github.com/aslakjohansen/sdu-sest-student-visit}
%    \end{center}
%\end{frame}

%\subsection{Tekst Editor}
%\begin{frame}
%    \frametitle{Tekst Editor}
%    Vi skal skrive noget kode. Det gøres i en \textsl{tekst editor}.
%    
%    \vspace{1cm}
%    \begin{center}
%        \begin{tabular}{l|c|c|c|}
%            \rowcolor{blue!25} \cellcolor{blue!0}
%                                  & \textbf{Windows} & \textbf{OSX} & \textbf{Linux} \\
%            \hline
%            \rowcolor{blue!10} \cellcolor{blue!25}
%            \textbf{Installeret:} & Notepad          & TextEdit     & \ldots \\
%            \hline
%            \rowcolor{blue!10} \cellcolor{blue!25}
%            \textbf{Bedre valg:}  & Sublime Text     & Sublime Text & Gedit/Kate/Sublime Text \\
%            \hline
%        \end{tabular}
%    \end{center}
%    \vspace{1cm}
%    
%    Sublime Text download: \url{https://www.sublimetext.com/download}
%\end{frame}

%%%%%%%%%%%%%%%%%%%%%%%%%%%%%%%%%%%%%%%%%%%%%%%%%%%%%%%%%%%%%%%%
%%%%%%%%%%%%%%%%%%%%%%%%%%%%%%%%%%%%%%%%%%%%%%%% HTML Dokumenter
%\section{HTML Dokumenter}
%\begin{frame}
%    \vspace{25mm}
%    \begin{center}
%        \Huge{Part 1:\\HTML Dokumenter}
%    \end{center}
%\end{frame}

%\subsection{Arbejde med HTML Filer}
%\begin{frame}
%    \frametitle{Arbejde med HTML Filer}
%    \vspace{20mm}
%    \begin{center}
%      \scalebox{1.6}{\includeSVG{html_flow}}
%    \end{center}
%\end{frame}

%\subsection{Hello, World}
%\begin{frame}
%    \frametitle{Hello, World}
%    \inputminted{html}{../src/frontend/part1_html_hello/index.html}
%\end{frame}

%\subsubsection{Indhold}
%\begin{frame}
%    \frametitle{Hello, World - Indhold}
%    \inputminted[highlightlines={6-8}]{html}{../src/frontend/part1_html_hello/index.html}
%\end{frame}

%\subsubsection{Metadata}
%\begin{frame}
%    \frametitle{Hello, World - Metadata}
%    \inputminted[highlightlines={2-5}]{html}{../src/frontend/part1_html_hello/index.html}
%\end{frame}

%\subsubsection{Dokumentets Titel}
%\begin{frame}
%    \frametitle{Hello, World - Dokumentets Titel}
%    \inputminted[highlightlines={3}]{html}{../src/frontend/part1_html_hello/index.html}
%\end{frame}

%\subsubsection{Dokumentets Kodning}
%\begin{frame}
%    \frametitle{Hello, World - Dokumentets Kodning}
%    \inputminted[highlightlines={4}]{html}{../src/frontend/part1_html_hello/index.html}
%\end{frame}

%\subsection{Anatomien af et Tag}
%\begin{frame}
%    \frametitle{Anatomien af et Tag}
%    \vspace{3mm}
%    Overblik:
%    \begin{itemize}
%      \item Indkapsling: \texttt{<\tagname{div}> \ldots\  </\tagname{div}>}
%      \item Navn: \texttt{<\tagname{div}> \ldots\  </\tagname{div}>} $\rightarrow$ \texttt{\tagname{div}}
%      \item Shorthand hvis indhold er tomt: \texttt{<\tagname{div}></\tagname{div}>} $\rightarrow$ \texttt{<\tagname{div} />}
%      \item Parametre: \texttt{<\tagname{div} \tagkey{a}=\tagvalue{"1"} \tagkey{b}=\tagvalue{"2"} />}
%      \item Identitet: \texttt{<\tagname{div} \tagkey{id}=\tagvalue{"data"} />}
%    \end{itemize}
%    
%    \vspace{5mm}
%    Dokumentet organiseres ved at
%    \begin{itemize}
%      \item placere tags indeni hinanden.
%      \item placere tags efter hinancen, eventuelt med tekst imellem.
%    \end{itemize}
%\end{frame}

%\subsection{Eksempel}
%\begin{frame}
%    \frametitle{Eksempel}
%    \inputminted[breaklines=true]{html}{../src/frontend/part1_html/index.html}
%\end{frame}

%%%%%%%%%%%%%%%%%%%%%%%%%%%%%%%%%%%%%%%%%%%%%%%%%%%%%%%%%%%%%%%%
%%%%%%%%%%%%%%%%%%%%%%%%%%%%%%%%%%%%%%%%%%% Logik med JavaScript
%\section{Logik med JavaScript}
%\begin{frame}
%    \vspace{25mm}
%    \begin{center}
%        \Huge{Part 2:\\Logik med JavaScript}
%    \end{center}
%\end{frame}

%\subsection{JavaScript?}
%\begin{frame}
%    \frametitle{JavaScript?}
%    Et programmeringssprog der
%    \begin{itemize}
%      \item primært er designet til at kunne afvikles i en browser.
%      \item kan manipulere den HTML datastruktur der vises i browseren.
%      \item kan reagere på handlinger på hjemmesiden.
%      \item kan interagere med omverdenen.
%      \item kan udvides igennem moduler.
%    \end{itemize}
%\end{frame}

%\subsection{Arbejde med Javascript Filer}
%\begin{frame}
%    \frametitle{Arbejde med Javascript Filer}
%    \vspace{7mm}
%    \begin{center}
%      \scalebox{1.6}{\includeSVG{js_flow}}
%    \end{center}
%\end{frame}

%\subsubsection{Dokumentet}

%\begin{frame}
%    \frametitle{Dokument}
%    \inputminted{html}{../src/frontend/part2_hello/index.html}
%\end{frame}

%\begin{frame}
%    \frametitle{Dokumentet}
%    \inputminted[highlightlines={11}]{html}{../src/frontend/part2_hello/index.html}
%\end{frame}

%\begin{frame}
%    \frametitle{Dokumentet}
%    \inputminted[highlightlines={9}]{html}{../src/frontend/part2_hello/index.html}
%\end{frame}

%\subsection{Udvikling}

%\subsubsection{Firefox}
%\begin{frame}
%    \frametitle{Udvikling - Firefox}
%    \begin{center}
%      \includeBitmap{firefox-devtools.png}{13cm}
%    \end{center}
%\end{frame}

%\subsubsection{Chrome og Chromium}
%\begin{frame}
%    \frametitle{Udvikling - Chrome og Chromium}
%    \begin{center}
%      \includeBitmap{chromium-devtools.png}{13cm}
%    \end{center}
%\end{frame}

%\subsection{Konsollen}
%\begin{frame}
%    \frametitle{Konsollen}
%    Vores første JavaScript program:
%    \inputminted{javascript}{../src/frontend/part2_hello/logic.js}
%\end{frame}

%\subsection{Logik}
%\begin{frame}
%    \frametitle{Logik}
%    \inputminted{javascript}{../src/frontend/part2_console/logic.js}
%\end{frame}
%\begin{frame}
%    \frametitle{Logik}
%    \inputminted[highlightlines={4,5,6}]{javascript}{../src/frontend/part2_console/logic.js}
%\end{frame}
%\begin{frame}
%    \frametitle{Logik}
%    \inputminted[highlightlines={7}]{javascript}{../src/frontend/part2_console/logic.js}
%\end{frame}
%\begin{frame}
%    \frametitle{Logik}
%    \inputminted[highlightlines={5,8}]{javascript}{../src/frontend/part2_console/logic.js}
%\end{frame}

%\subsection{Manipulation med HTML}
%\begin{frame}
%    \frametitle{Manipulation med HTML}
%    \vspace{-1mm}
%    \inputminted[highlightlines={},breaklines=true]{javascript}{../src/frontend/part2_html/logic.js}
%\end{frame}
%\begin{frame}
%    \frametitle{Manipulation med HTML}
%    \vspace{-1mm}
%    \inputminted[highlightlines={1,9,11,13},breaklines=true]{javascript}{../src/frontend/part2_html/logic.js}
%\end{frame}

%%%%%%%%%%%%%%%%%%%%%%%%%%%%%%%%%%%%%%%%%%%%%%%%%%%%%%%%%%%%%%%%
%%%%%%%%%%%%%%%%%%%%%%%%%%%%%%%%%%%%%%%%%%%%%%%%%%% Tekststrenge
%\section{Strukturerede Data}
%\begin{frame}
%    \vspace{25mm}
%    \begin{center}
%        \Huge{Part 3:\\Strukturerede Data}
%    \end{center}
%\end{frame}

%\subsection{Datastrukturer}
%\subsubsection{JavaScript Objekter}

%\begin{frame}
%    \frametitle{Datastrukturer - JavaScript Objekter}
%    \vspace{-1mm}
%    \begin{center}
%      \scalebox{1.3}{\includeSVG{object-struct0}}
%    \end{center}
%\end{frame}
%\begin{frame}
%    \frametitle{Datastrukturer - JavaScript Objekter}
%    \vspace{-1mm}
%    \begin{center}
%      \scalebox{1.3}{\includeSVG{object-struct1}}
%    \end{center}
%\end{frame}
%\begin{frame}
%    \frametitle{Datastrukturer - JavaScript Objekter}
%    \vspace{-1mm}
%    \begin{center}
%      \scalebox{1.3}{\includeSVG{object-struct2}}
%    \end{center}
%\end{frame}

%\begin{frame}
%    \frametitle{Datastrukturer - JavaScript Objekter, Anvendelse af}
%    \vspace{-1mm}
%    \inputminted[fontsize=\footnotesize]{javascript}{../src/frontend/part3_object/logic.js}
%\end{frame}
%\begin{frame}
%    \frametitle{Datastrukturer - JavaScript Objekter, Anvendelse af}
%    \vspace{-1mm}
%    \inputminted[fontsize=\footnotesize,highlightlines={1}]{javascript}{../src/frontend/part3_object/logic.js}
%\end{frame}
%\begin{frame}
%    \frametitle{Datastrukturer - JavaScript Objekter, Anvendelse af}
%    \vspace{-1mm}
%    \inputminted[fontsize=\footnotesize,highlightlines={4,8}]{javascript}{../src/frontend/part3_object/logic.js}
%\end{frame}
%\begin{frame}
%    \frametitle{Datastrukturer - JavaScript Objekter, Anvendelse af}
%    \vspace{-1mm}
%    \inputminted[fontsize=\footnotesize,highlightlines={9}]{javascript}{../src/frontend/part3_object/logic.js}
%\end{frame}
%\begin{frame}
%    \frametitle{Datastrukturer - JavaScript Objekter, Anvendelse af}
%    \vspace{-1mm}
%    \inputminted[fontsize=\footnotesize]{javascript}{../src/frontend/part3_object/logic.js}
%\end{frame}

%\subsection{Tekststrenge}
%\begin{frame}
%    \frametitle{Tekststrenge}
%    \vspace{3mm}
%    En tekststreng (eller bare \textsl{string}) er en sekvens af tegn.
%    
%    \vspace{5mm}
%    Vi har før set denne:
%    \vspace{-1mm}
%    \begin{center}
%      \scalebox{1.6}{\includeSVG{string-hello}}
%    \end{center}
%    
%    Men vi kan også oprette denne:
%    \vspace{-1mm}
%    \begin{center}
%      \scalebox{1.6}{\includeSVG{string-struct}}
%    \end{center}
%\end{frame}

%\subsection{JSON Formatet}
%\begin{frame}
%    \frametitle{JSON Formatet}
%    \vspace{-1mm}
%    \inputminted[fontsize=\footnotesize]{javascript}{../src/frontend/part3_object/logic.js}
%\end{frame}
%\begin{frame}
%    \frametitle{JSON Formatet}
%    \vspace{-1mm}
%    \inputminted[fontsize=\footnotesize,highlightlines={1}]{javascript}{../src/frontend/part3_json/logic.js}
%    
%    \pause
%    \vspace{1mm}
%    JSON er en måde at konvertere en datastruktur til en streng, og tilbage igen.
%\end{frame}

%%%%%%%%%%%%%%%%%%%%%%%%%%%%%%%%%%%%%%%%%%%%%%%%%%%%%%%%%%%%%%%%
%%%%%%%%%%%%%%%%%%%%%%%%%%%%%%%%%%%%%%%%%%%%%% Publish Subscribe
%\section{Publish Subscribe}
%\begin{frame}
%    \vspace{25mm}
%    \begin{center}
%        \Huge{Part 4:\\Data med Publish Subscribe}
%    \end{center}
%\end{frame}

%\subsection{Overblik}
%\begin{frame}
%    \frametitle{Overblik}
%    \vspace{10mm}
%    \begin{center}
%      \scalebox{1.6}{\includeSVG{pubsub_flow}}
%    \end{center}
%\end{frame}
%\begin{frame}
%    \frametitle{Overblik}
%    \vspace{10mm}
%    \begin{center}
%      \scalebox{1.6}{\includeSVG{pubsub_flow0}}
%    \end{center}
%\end{frame}

%\subsection{Modtagelse af Beskeder}
%\subsubsection{HTML Fil}
%\begin{frame}
%    \frametitle{Modtagelse af Beskeder - HTML Fil}
%    \vspace{-1mm}
%    \inputminted[fontsize=\footnotesize,breaklines=true,highlightlines={11}]{html}{../src/frontend/part4_recv/index.html}
%\end{frame}
%\subsubsection{JavaScript Fil}
%\begin{frame}
%    \frametitle{Modtagelse af Beskeder - JavaScript Fil}
%    \vspace{-1mm}
%    \inputminted[fontsize=\footnotesize,highlightlines={}]{javascript}{../src/frontend/part4_recv/logic.js}
%\end{frame}
%\begin{frame}
%    \frametitle{Modtagelse af Beskeder - JavaScript Fil}
%    \vspace{-1mm}
%    \inputminted[fontsize=\footnotesize,highlightlines={1}]{javascript}{../src/frontend/part4_recv/logic.js}
%\end{frame}
%\begin{frame}
%    \frametitle{Modtagelse af Beskeder - JavaScript Fil}
%    \vspace{-1mm}
%    \inputminted[fontsize=\footnotesize,highlightlines={2}]{javascript}{../src/frontend/part4_recv/logic.js}
%\end{frame}
%\begin{frame}
%    \frametitle{Modtagelse af Beskeder - JavaScript Fil}
%    \vspace{-1mm}
%    \inputminted[fontsize=\footnotesize,highlightlines={5}]{javascript}{../src/frontend/part4_recv/logic.js}
%\end{frame}
%\begin{frame}
%    \frametitle{Modtagelse af Beskeder - JavaScript Fil}
%    \vspace{-1mm}
%    \inputminted[fontsize=\footnotesize,highlightlines={6,13}]{javascript}{../src/frontend/part4_recv/logic.js}
%\end{frame}
%\begin{frame}
%    \frametitle{Modtagelse af Beskeder - JavaScript Fil}
%    \vspace{-1mm}
%    \inputminted[fontsize=\footnotesize,highlightlines={7,9}]{javascript}{../src/frontend/part4_recv/logic.js}
%\end{frame}
%\begin{frame}
%    \frametitle{Modtagelse af Beskeder - JavaScript Fil}
%    \vspace{-1mm}
%    \inputminted[fontsize=\footnotesize,highlightlines={10}]{javascript}{../src/frontend/part4_recv/logic.js}
%\end{frame}

%\subsection{Behandling af Beskeder}
%\begin{frame}
%    \frametitle{Behandling af Beskeder}
%    \vspace{-1mm}
%    \inputminted[fontsize=\tiny,highlightlines={}]{javascript}{../src/frontend/part4_dashboard/logic.js}
%\end{frame}
%\begin{frame}
%    \frametitle{Behandling af Beskeder}
%    \vspace{-1mm}
%    \inputminted[fontsize=\tiny,highlightlines={14}]{javascript}{../src/frontend/part4_dashboard/logic.js}
%\end{frame}
%\begin{frame}
%    \frametitle{Behandling af Beskeder}
%    \vspace{-1mm}
%    \inputminted[fontsize=\tiny,highlightlines={15,19,20,21,23}]{javascript}{../src/frontend/part4_dashboard/logic.js}
%\end{frame}
%\begin{frame}
%    \frametitle{Behandling af Beskeder}
%    \vspace{-1mm}
%    \inputminted[fontsize=\tiny,highlightlines={17,18,22,24}]{javascript}{../src/frontend/part4_dashboard/logic.js}
%\end{frame}

%%%%%%%%%%%%%%%%%%%%%%%%%%%%%%%%%%%%%%%%%%%%%%%%%%%%%%%%%%%%%%%%
%%%%%%%%%%%%%%%%%%%%%%%%%%%%%%%%%%%%%%%%%%%%%%%%%%%%%%%%% Farver
%\section{Farver}
%\begin{frame}
%    \vspace{25mm}
%    \begin{center}
%        \Huge{Part 5:\\Farver}
%    \end{center}
%\end{frame}

%\subsection{RGB Farvekuben}
%\begin{frame}
%    \frametitle{RGB Farvekuben}
%    \vspace{-11mm}
%    \begin{center}
%      \scalebox{1.6}{\includeSVG{colors}}
%    \end{center}
%\end{frame}

%\subsection{Hexadecimal Repræsentation}
%\begin{frame}
%    \frametitle{Hexadecimal Repræsentation}
%    \vspace{3mm}
%    Altså, en farve i en computer er et punkt i et 3d koordinatsystem med en grundfarve på hver akse, og hvor hver akse spænder mellem 0 og 1.
%    
%    \vspace{5mm}
%    I praksis arbejder man med 8 bit per \textsl{kanal}. Dette betyder at hver akse (aka kanal) inddeles i $2^8=256$ niveauer. Dette giver $256^3=16777216$ forskellige farver.
%    
%    \vspace{5mm}
%    I HTML filer beskrives disse værdier typisk hexadecimalt, altså ved hjælp af et 16-tals system. 
%    
%    \vspace{5mm}
%    Der følges formatet "\texttt{rrggbb}". I hver farves kodning repræsenterer det højre tegn en værdi mellem 0 og 15 (herunder ${a\mapsto 10, b \mapsto 11, c \mapsto 12, d \mapsto 13, e \mapsto 14, f \mapsto 15}$), og det venstre tegn er 16 gange så meget værd.
%    
%     \vspace{5mm}
%     \textcolor{black}{\texttt{000000}}
%     \textcolor{red}{\texttt{ff0000}}
%     \textcolor{green}{\texttt{00ff00}}
%     \textcolor{blue}{\texttt{0000ff}}
%     \textcolor{yellow}{\texttt{ffff00}}
%     \textcolor{cyan}{\texttt{00ffff}}
%     \textcolor{magenta}{\texttt{ff00ff}}
%\end{frame}

%%\subsection{Repræsentation med Tripletter}
%%\begin{frame}
%%    \frametitle{Repræsentation med Tripletter}
%%\end{frame}

%\subsection{Tærskelværdier}
%\begin{frame}
%    \frametitle{Tærskelværdier}
%    \vspace{-2mm}
%    \inputminted[fontsize=\tiny,highlightlines={}]{javascript}{../src/frontend/part5_threshold/logic.js}
%\end{frame}
%\begin{frame}
%    \frametitle{Tærskelværdier}
%    \vspace{-2mm}
%    \inputminted[fontsize=\tiny,highlightlines={9-15,26}]{javascript}{../src/frontend/part5_threshold/logic.js}
%\end{frame}

%%%%%%%%%%%%%%%%%%%%%%%%%%%%%%%%%%%%%%%%%%%%%%%%%%%%%%%%%%%%%%%%
%%%%%%%%%%%%%%%%%%%%%%%%%%%%%%%%%%%%%%%%%%%%%%%%%%%%%%%%%% Plots
%\section{Plots}
%\begin{frame}
%    \vspace{25mm}
%    \begin{center}
%        \Huge{Part 6:\\Plots}
%    \end{center}
%\end{frame}

%\subsection{HTML Fil}
%\begin{frame}
%    \frametitle{HTML Fil}
%    \vspace{-1mm}
%    \inputminted[fontsize=\footnotesize,breaklines=true,highlightlines={10,13}]{html}{../src/frontend/part6_plot/index.html}
%\end{frame}

%\subsection{JavaScript Fil}
%\begin{frame}
%    \frametitle{JavaScript Fil}
%    \vspace{-2mm}
%    \begin{multicols}{2}
%      \inputminted[fontsize=\tiny,highlightlines={}]{javascript}{../src/frontend/part6_plot/logic.js}
%    \end{multicols}
%\end{frame}
%\begin{frame}
%    \frametitle{JavaScript Fil}
%    \vspace{-2mm}
%    \begin{multicols}{2}
%      \inputminted[fontsize=\tiny,highlightlines={1-3,39-40,44-45}]{javascript}{../src/frontend/part6_plot/logic.js}
%    \end{multicols}
%\end{frame}
%\begin{frame}
%    \frametitle{JavaScript Fil}
%    \vspace{-2mm}
%    \begin{multicols}{2}
%      \inputminted[fontsize=\tiny,highlightlines={42}]{javascript}{../src/frontend/part6_plot/logic.js}
%    \end{multicols}
%\end{frame}
%\begin{frame}
%    \frametitle{JavaScript Fil}
%    \vspace{-2mm}
%    \begin{multicols}{2}
%      \inputminted[fontsize=\tiny,highlightlines={43-46}]{javascript}{../src/frontend/part6_plot/logic.js}
%    \end{multicols}
%\end{frame}

%%%%%%%%%%%%%%%%%%%%%%%%%%%%%%%%%%%%%%%%%%%%%%%%%%%%%%%%%%%%%%%%
%%%%%%%%%%%%%%%%%%%%%%%%%%%%%%%%%%%%%%%%%%%%%%%%%%%%%%%%%%% Test
%\section{Test}
%\begin{frame}
%    \vspace{25mm}
%    \begin{center}
%        \Huge{Part 7:\\Test med Manuel Publicering}
%    \end{center}
%\end{frame}

%\subsection{Manuel Publicering}
%\subsubsection{HTML Fil}
%\begin{frame}
%    \frametitle{Modtagelse af Beskeder - HTML Fil}
%    \vspace{-2mm}
%    \inputminted[fontsize=\footnotesize,breaklines=true,highlightlines={}]{html}{../src/frontend/part7_publish/index.html}
%\end{frame}
%\begin{frame}
%    \frametitle{Modtagelse af Beskeder - HTML Fil}
%    \vspace{-2mm}
%    \inputminted[fontsize=\footnotesize,breaklines=true,highlightlines={9,11,12}]{html}{../src/frontend/part7_publish/index.html}
%\end{frame}

%\subsubsection{JavaScript Fil}
%\begin{frame}
%    \frametitle{Modtagelse af Beskeder - JavaScript Fil}
%    \vspace{-1mm}
%    \inputminted[fontsize=\normalsize,highlightlines={}]{javascript}{../src/frontend/part7_publish/logic.js}
%\end{frame}
%\begin{frame}
%    \frametitle{Modtagelse af Beskeder - JavaScript Fil}
%    \vspace{-1mm}
%    \inputminted[fontsize=\normalsize,highlightlines={1}]{javascript}{../src/frontend/part7_publish/logic.js}
%\end{frame}
%\begin{frame}
%    \frametitle{Modtagelse af Beskeder - JavaScript Fil}
%    \vspace{-1mm}
%    \inputminted[fontsize=\normalsize,highlightlines={3,4}]{javascript}{../src/frontend/part7_publish/logic.js}
%\end{frame}
%\begin{frame}
%    \frametitle{Modtagelse af Beskeder - JavaScript Fil}
%    \vspace{-1mm}
%    \inputminted[fontsize=\normalsize,highlightlines={14,6}]{javascript}{../src/frontend/part7_publish/logic.js}
%\end{frame}
%\begin{frame}
%    \frametitle{Modtagelse af Beskeder - JavaScript Fil}
%    \vspace{-1mm}
%    \inputminted[fontsize=\normalsize,highlightlines={9}]{javascript}{../src/frontend/part7_publish/logic.js}
%\end{frame}

%\subsection{Robusthed}
%\begin{frame}
%    \frametitle{Robusthed}
%    Hvad sker der hvis vi sender strengen \say{\textcolor{purple}{\texttt{\{"temperature":}}} i stedet for et gyldigt JSON objekt?
%    \pause
%    
%    \vspace{5mm}
%    Vores modtager holder op med at reagere!
%    \pause
%    
%    \vspace{5mm}
%    Et program siges at være \textcolor{teal}{\textsl{robust}} hvis det reagerer hensigtsmæssigt (e.g., det ikke crasher) når det gives fejlbehæftede input.
%\end{frame}
%\begin{frame}
%    \frametitle{JavaScript Fil}
%    \vspace{-2mm}
%    \begin{multicols}{2}
%      \inputminted[fontsize=\tiny,highlightlines={31-35}]{javascript}{../src/frontend/part7_plot/logic.js}
%    \end{multicols}
%\end{frame}

%%%%%%%%%%%%%%%%%%%%%%%%%%%%%%%%%%%%%%%%%%%%%%%%%%%%%%%%%%%%%%%%
%%%%%%%%%%%%%%%%%%%%%%%%%%%%%%%%%%%%%%%%%%%%%%%%%%%% Chat System
%\section{Chat System}
%\begin{frame}
%    \vspace{25mm}
%    \begin{center}
%        \Huge{Part 8:\\Chat System}
%    \end{center}
%\end{frame}

%\subsection{Øvelse}
%\begin{frame}
%    \frametitle{Øvelse}
%    \begin{itemize}
%      \item Vælg et nyt topic navn.
%      \item Lav en HTML side med en knap og et tekstfelt. Når knappen aktiveres skal indholdet af tekstfeltet publiceres til det valgte topic.
%      \item Opret en subscription til det valgte topic. Når data ankommer på dette topic vises det et sted på siden.
%    \end{itemize}
%\end{frame}

%%%%%%%%%%%%%%%%%%%%%%%%%%%%%%%%%%%%%%%%%%%%%%%%%%%%%%%%%%%%%%%
%%%%%%%%%%%%%%%%%%%%%%%%%%%%%%%%%%%%%%%%%%%%%%%%%%%%% Questions
\section{Spørgsmål}
\begin{frame}
    \frametitle{\textbf{Spørgsmål?}}
    \vspace{-15mm}
    \begin{center}
        \includegraphics[scale=0.4]{./figs/Boy-asking-question.pdf}
    \end{center}
    \vspace{-25mm}
    \scalebox{0.2}{\url{https://openclipart.org/detail/238687/boy-thinking-of-question}}
\end{frame}

